% Budget: ca. 0,1 Seiten. Verantwortlich: Person C
% Forschungsfrage explizit beantworten (Qualität und Produktivität), dann Ausblick.
\section{Fazit}
\label{sec:fazit}

Aus der aktuellen Forschung lassen sich fünf Gestaltungsprinzipien ableiten: Planen vor Ausführen, Autonomie nach Aufgabe und Risiko staffeln, knapper und aufgabenbezogener Kontext, mehrstufige automatische Verifikation sowie verbindliches Review kleiner Änderungen. Für die Codequalität sind Verifikation und Review am besten belegt, denn bestandene Tests allein garantieren keine Korrektheit, und mehr Autonomie erhöht das Sicherheitsrisiko. Die Produktivität steigt vor allem bei unbekannten Codebasen und klar umrissenen Aufgaben, kann bei erfahrenen Entwicklern in vertrauten Großprojekten aber sinken. Nicht der Grad der Automatisierung entscheidet daher über den Nutzen, sondern wie gut der Workflow Autonomie, Kontext und Prüfung aufeinander abstimmt.

Künftige Arbeiten sollten diese Prinzipien in kontrollierten Studien mit realen Teams vergleichen und ihre langfristige Wirkung auf Wartbarkeit und Sicherheit untersuchen.
